
\section{Literature Survey}
This survey covers every work the project builds on: the foundations on Transformers, RAG, hallucination
and long-context evaluation that motivated the memory layer, and the retrieval, memory, grounding and
benchmark literature that each stage of the implemented layer draws on. Entries are grouped by theme and
each ends with the limitation that motivated a design decision in Chapter~3.

\subsection{Transformer Architecture: Attention Is All You Need (Vaswani et al.\ 2017)}
Vaswani et al.\ \cite{vaswani2017} introduced the Transformer, which replaces recurrence with
self-attention so that every token can attend to every other token in parallel. It is the architecture
of every LLM used in this project, and the reason the context window exists at all: attention is computed
over a fixed-length window of tokens. \\
\textbf{Limitation:} attention costs quadratic time and memory in the sequence length, so simply making
the window longer is expensive, and nothing outside the window is visible to the model.

\subsection{Retrieval-Augmented Generation for Knowledge-Intensive NLP (Lewis et al.\ 2020)}
Lewis et al.\ \cite{lewis2020} combined a dense retriever with a sequence-to-sequence generator so that
the model conditions on retrieved passages instead of relying only on parametric knowledge, improving
factual accuracy on open-domain QA. RAG is the starting point of our Phase-1 design (vector RAG baseline). \\
\textbf{Limitation:} RAG is stateless across interactions and has no notion of the user, of time, or of
which retrieved passage supersedes another.

\subsection{Retrieval-Augmented Generation for LLMs: A Survey (Gao et al.\ 2023)}
Gao et al.\ \cite{gao2023} survey naive, advanced and modular RAG, covering chunking, embedding, hybrid
retrieval, re-ranking and query rewriting, and report consistent reductions in hallucination from
grounding. Several of the ``advanced RAG'' components they list (hybrid search, re-ranking, query
expansion) are implemented in our read path. \\
\textbf{Limitation:} the surveyed systems index static documents; conversational memory with updates and
personal facts is not addressed.

\subsection{A Survey on Hallucination in Large Language Models (Huang et al.\ 2024)}
Huang et al.\ \cite{huang2024} classify hallucinations into factuality and faithfulness errors, trace their
causes to data, training and inference, and review detection and mitigation methods including retrieval
augmentation. \\
\textbf{Limitation:} the survey shows that retrieval reduces but does not remove hallucination, and that
detection methods trade precision for recall, which is exactly what Phase 1 measured on RAGTruth.

\subsection{Lost in the Middle: How Language Models Use Long Contexts (Liu et al.\ 2023)}
Liu et al.\ \cite{liu2023} showed that when the relevant document is placed in the middle of a long
prompt, accuracy drops sharply compared with placing it at the beginning or end, producing a U-shaped
curve across positions. \\
\textbf{Limitation:} models do not use their context uniformly, so a memory that is ``in the window'' is
not necessarily used; this motivates injecting a few notes rather than the whole history.

\subsection{NoLiMa: Long-Context Evaluation Beyond Literal Matching (Modarressi et al.\ 2025)}
NoLiMa \cite{nolima2025} removes lexical overlap between the question and the needle so that the model
must reason associatively, and finds that accuracy at 32k tokens falls far below the short-context
accuracy of the same models. \\
\textbf{Limitation:} long-context performance claims based on literal needle matching are optimistic; our
benchmark therefore paraphrases every question relative to its stored evidence.

\subsection{Why Does the Effective Context Length of LLMs Fall Short? (An et al.\ 2025)}
An et al.\ \cite{an2025} attribute the gap between the trained and the usable context length to the
left-skewed distribution of relative positions seen in training, and propose position remapping (STRING)
to recover performance without retraining. \\
\textbf{Limitation:} even with the fix, the effective context is a fraction of the nominal window, so
history that exceeds it must be selected rather than appended.

\subsection{Needle-in-a-Haystack Evaluation for Long-Context LLMs (Kamradt 2023; Nelson 2025)}
The needle-in-a-haystack protocol \cite{niah2025} inserts one fact at a chosen depth in a long filler
document and measures whether the model can retrieve it; results are reported as depth-by-length heatmaps
that show degradation with length and position. Phase 1 reported such heatmaps from the literature; Phase
2 reproduces the protocol on our own model (Section~\ref{sec:niah}). \\
\textbf{Limitation:} a single needle in unrelated filler is easier than a real history that contains many
similar facts, updates and distractors; the protocol measures length, not interference.

\subsection{Dense Retrieval and Retrieval-Augmented Pre-training}
\textbf{DPR (Karpukhin et al.\ 2020)} \cite{dpr} trains a bi-encoder so that questions and passages are
compared by a single inner product, and showed that learned dense retrieval beats BM25 on open-domain QA
when the question and passage share little vocabulary. This is the justification for the dense channel of
our first stage. \\
\textbf{Limitation:} bi-encoders miss exact matches of rare strings (names, model numbers), which is why a
lexical channel is still needed.

\textbf{REALM (Guu et al.\ 2020)} \cite{realm} and \textbf{RETRO (Borgeaud et al.\ 2022)} \cite{retro}
integrate retrieval into pre-training itself: REALM back-propagates through a latent retriever, and RETRO
conditions a 7B model on chunks retrieved from a two-trillion-token database, matching models 25 times
larger. They demonstrate that retrieving a few relevant chunks is worth more than scaling parameters. \\
\textbf{Limitation:} both require retraining the language model; a memory layer must work with a frozen,
off-the-shelf model.

\textbf{Fusion-in-Decoder (Izacard and Grave 2021)} \cite{fid} encodes each retrieved passage
independently and lets the decoder attend over all of them, and found that accuracy keeps improving up to
about 100 passages while a single long concatenation does not. \\
\textbf{Limitation:} needs an encoder--decoder model; with a decoder-only LLM the practical substitute is
to inject a small number of high-precision notes, which is what the token budget in our layer enforces.

\subsection{Lexical Retrieval, Late Interaction and Fusion}
\textbf{BM25 (Robertson and Zaragoza 2009)} \cite{bm25} remains the standard lexical ranking function and is
the sparse channel of our retriever. \textbf{ColBERT (Khattab and Zaharia 2020)} \cite{colbert} keeps one
vector per token and scores by late interaction (MaxSim), getting close to cross-encoder quality at
bi-encoder cost. \textbf{Reciprocal Rank Fusion (Cormack et al.\ 2009)} \cite{rrf} merges ranked lists from
several retrievers using only ranks, which makes it robust to incomparable score scales. \\
\textbf{Limitation:} RRF with equal weights lets a weak channel inject noise, and BM25 returns near-zero
``matches'' on shared frequent words; we use a weighted RRF and a score floor on the BM25 channel
(Section~\ref{sec:fusion}). ColBERT indexes are an order of magnitude larger than single-vector indexes,
so we kept the simpler bi-encoder.

\textbf{HyDE (Gao et al.\ 2023)} \cite{hyde} performs zero-shot retrieval by asking an LLM to write a
hypothetical answer document and embedding that instead of the query, closing the vocabulary gap between
questions and passages. \\
\textbf{Limitation:} one LLM call per query; our rule-based query expansion is a zero-cost approximation
of the same idea for the lexical channel.

\textbf{Sentence-BERT (Reimers and Gurevych 2019)} \cite{sbert} produces dense sentence embeddings suited to
semantic similarity search, and \textbf{FAISS (Johnson et al.\ 2021)} \cite{faiss} provides fast exact and
approximate nearest-neighbour search over them; both are used directly in the layer.

\subsection{Cross-Encoder Re-ranking}
\textbf{Nogueira and Cho (2019)} \cite{nogueira2019} showed that a BERT cross-encoder that reads the query
and the passage jointly re-ranks first-stage candidates far more accurately than bi-encoders. We use the
MS~MARCO MiniLM cross-encoder as the re-ranking stage. \\
\textbf{Limitation:} cross-encoders are trained on web passages; on short, first-person conversational
memories they are unreliable unless the memory text is normalised (Section~\ref{sec:normalisation}), and
even then they are blended with first-stage evidence rather than trusted alone.

\subsection{Memory-Augmented Transformers}
\textbf{Memorizing Transformers (Wu et al.\ 2022)} \cite{memtrans} add a $k$NN look-up over the keys and
values of past tokens to one attention layer, extending the usable context to hundreds of thousands of
tokens without retraining from scratch. \textbf{LongMem (Wang et al.\ 2023)} \cite{longmem} decouples the
memory into a frozen backbone that writes key--value pairs and a trainable side-network that reads them. \\
\textbf{Limitation:} both store opaque activations tied to one model; the memory cannot be inspected,
edited, or moved to another model. Our layer stores natural-language notes, which are model-independent
and auditable.

\subsection{Memory for LLM Agents}
\textbf{Generative Agents (Park et al.\ 2023)} \cite{park2023} gave simulated agents a \emph{memory stream}
of natural-language observations and retrieved memories with a weighted sum of \emph{recency}
(exponential decay), \emph{importance} (rated once at write time) and \emph{relevance} (embedding
similarity). This triple is the basis of our memory scoring stage. \\
\textbf{Limitation:} the three scores are added, so a recent but irrelevant memory can outrank a relevant
old one; we found exactly this failure mode in our first implementation and changed to a relevance-gated,
multiplicative combination (Section~\ref{sec:scoring}).

\textbf{Reflexion (Shinn et al.\ 2023)} \cite{reflexion} stores verbal self-reflections in an episodic
memory buffer so that an agent improves across trials. \textbf{MemGPT (Packer et al.\ 2023)}
\cite{packer2023} treats the context window as main memory and an external store as disk, letting the LLM
page information in and out through function calls. \textbf{MemoryBank (Zhong et al.\ 2024)}
\cite{zhong2024} adds an Ebbinghaus-style forgetting curve so that memories that are not revisited fade. \\
\textbf{Limitation:} all three rely on the LLM itself to decide what to store and recall, which is
expensive and hard to evaluate; our layer keeps storage and retrieval deterministic and model-independent.

\textbf{Mem0 (Chhikara et al.\ 2025)} \cite{mem0} extracts salient facts from each turn, stores them in a
vector database and resolves each new fact against existing ones with \emph{add / update / delete / no-op}
operations; a graph variant stores entities and relations. It reports large token savings over
full-context prompting on LoCoMo. \\
\textbf{Limitation:} fact extraction and conflict resolution are performed by an LLM call per turn. We
approximate the update operation with an embedding-similarity plus update-cue rule that costs nothing at
write time.

\textbf{Survey of memory mechanisms (Zhang et al.\ 2024)} \cite{memsurvey} organises agent memory along
the axes of \emph{what} is stored (raw text, summaries, structured facts), \emph{how} it is managed
(write, update, forget) and \emph{how} it is read (relevance, recency, importance). Our design covers all
three axes with explicit, switchable components so that each can be ablated.

\subsection{Graph-Based Retrieval}
\textbf{GraphRAG (Edge et al.\ 2024)} \cite{graphrag} extracts an entity--relation graph from a corpus and
answers questions by combining vector search with graph traversal, which helps questions whose evidence is
spread over several documents. Our entity graph is a light-weight version: memories are linked through
shared named entities and retrieval expands one hop from the top candidates. \\
\textbf{Limitation:} GraphRAG builds the graph with many LLM calls; our rule-based extraction is free but
only captures explicit surface entities.

\subsection{Long Context: Distraction, Retrieval versus Length, RULER}
\textbf{Shi et al.\ (2023)} \cite{distracted} showed that adding a single irrelevant but topically similar
sentence to a reasoning problem sharply lowers LLM accuracy, i.e.\ models are easily distracted by
irrelevant context. \textbf{Xu et al.\ (2024)} \cite{retrievalmeetslong} compared a 4k-context model with
retrieval against 16k/32k long-context models and found that retrieval of a few chunks matches or beats
the long-context model at a fraction of the cost. \textbf{RULER (Hsieh et al.\ 2024)} \cite{ruler}
generalises the needle test to multi-needle, multi-hop tracing and aggregation tasks and found that most
models' effective context is far below their claimed window. \\
\textbf{Limitation:} these studies measure the model; they do not say how to \emph{select} the few chunks
in a multi-session conversation. Our distractor facts and NIAH-style scaling experiment transfer their
methodology to the memory setting.

\subsection{Hallucination Corpora and Grounding Verification}
\textbf{RAGTruth (Niu et al.\ 2024)} \cite{ragtruth} is the word-level hallucination corpus for RAG outputs
that Phase 1 used; its baseline detectors (prompted GPT-3.5/4, SelfCheckGPT, LMvLM) reach high recall but
low precision. \textbf{SelfCheckGPT (Manakul et al.\ 2023)} \cite{selfcheck} detects hallucination by
sampling several answers and measuring their mutual inconsistency, and \textbf{FActScore (Min et al.\ 2023)}
\cite{factscore} breaks a generation into atomic facts and checks each against a knowledge source.
\textbf{Self-RAG (Asai et al.\ 2024)} \cite{selfrag} trains the model to emit reflection tokens that decide
whether to retrieve and whether the output is supported by the retrieved passage. \\
\textbf{Limitation:} SelfCheckGPT needs several extra generations and Self-RAG needs fine-tuning;
FActScore needs an external knowledge base. Our grounding gate borrows the atomic-fact attribution idea of
FActScore and the ``is it supported'' judgement of Self-RAG, but implements both as deterministic checks
against the injected notes (attributable, factual, same subject), which cost no extra model calls.

\subsection{Long-Term Conversational Memory Benchmarks}
\textbf{LoCoMo (Maharana et al.\ 2024)} \cite{locomo} and \textbf{LongMemEval (Wu et al.\ 2025)}
\cite{longmemeval} evaluate assistants on multi-session conversations with single-hop, multi-hop,
temporal-reasoning and abstention questions, and report that full-context prompting is both expensive and
surprisingly inaccurate, while retrieval over the history recovers most of the accuracy at a fraction of
the cost. \textbf{LongBench (Bai et al.\ 2024)} \cite{longbench} is the corresponding bilingual
multi-task suite for long documents. Our synthetic benchmark follows the LoCoMo/LongMemEval question
taxonomy. \\
\textbf{Limitation:} the published benchmarks need large API models to run and do not expose gold memory
identifiers; we built a smaller, fully controlled benchmark that a 1.5B-parameter local model can be
evaluated on reproducibly and that reports retrieval recall as well as answer accuracy.

\subsection{Small Instruction-Tuned Models (Qwen2.5)}
The Qwen2.5 family \cite{qwen25} provides instruction-tuned models from 0.5B to 72B parameters with 32k-token
context windows. We use the 1.5B model as the answering LLM for every method so that differences in the
results come from the memory strategy and not from the model.

\subsection{Summary}
Table~\ref{tab:litmap} maps each stage of the implemented memory layer to the work it draws on.
\begin{table}[htbp]\centering\small
\begin{tabular}{p{4.2cm}p{9.3cm}}\hline
\textbf{Stage of the memory layer} & \textbf{Source works} \\\hline
Dense channel (embeddings, FAISS) & DPR \cite{dpr}, Sentence-BERT \cite{sbert}, FAISS \cite{faiss} \\
Sparse channel, query expansion & BM25 \cite{bm25}, HyDE \cite{hyde} \\
Weighted rank fusion & RRF \cite{rrf} \\
Entity graph expansion & GraphRAG \cite{graphrag}, Mem0 graph memory \cite{mem0} \\
Cross-encoder re-ranking & Nogueira and Cho \cite{nogueira2019}, ColBERT \cite{colbert} \\
Memory scoring (recency, importance, relevance) & Generative Agents \cite{park2023}, MemoryBank \cite{zhong2024}, memory survey \cite{memsurvey} \\
Conflict-aware updates & Mem0 update operation \cite{mem0} \\
Token-budgeted injection & FiD \cite{fid}, retrieval vs.\ long context \cite{retrievalmeetslong}, distraction \cite{distracted} \\
Grounding gate & FActScore \cite{factscore}, Self-RAG \cite{selfrag}, RAGTruth \cite{ragtruth}, SelfCheckGPT \cite{selfcheck} \\
Benchmark design & LoCoMo \cite{locomo}, LongMemEval \cite{longmemeval}, RULER \cite{ruler}, NIAH \cite{niah2025} \\\hline
\end{tabular}
\caption{Mapping of the memory-layer stages to the literature.}\label{tab:litmap}
\end{table}

\section{Problem Statement and Objectives}
\noindent\textbf{Problem Statement 1:} A memory layer was proposed in Phase 1 but never built; without an
implementation its claimed benefits cannot be verified.\\
\textbf{Objective 1:} Implement the full pipeline as a reusable package with switchable stages.

\noindent\textbf{Problem Statement 2:} Plain vector retrieval cannot tell current facts from outdated ones and misses
facts connected only through entities.\\
\textbf{Objective 2:} Add conflict-aware recency and an entity graph, and measure their effect on temporal
and multi-hop questions.

\noindent\textbf{Problem Statement 3:} Hallucination detection in Phase 1 had precision of only 0.25.\\
\textbf{Objective 3:} Design a grounding gate that flags unsupported answers with high precision.

\noindent\textbf{Problem Statement 4:} Full-context prompting is the de-facto baseline but its cost and accuracy at
different history lengths were only taken from the literature in Phase 1.\\
\textbf{Objective 4:} Measure accuracy, tokens and latency of full-context versus memory-layer answering
on histories from 1k to 16k tokens with the same model.

\noindent\textbf{Problem Statement 5:} It is unknown which stages of a multi-stage memory pipeline actually matter.\\
\textbf{Objective 5:} Run a stage-by-stage ablation and a pruning study, and tune hyper-parameters on a
separate development split.
